\documentclass[11pt]{amsart}
\usepackage[T1]{fontenc}
\usepackage{lmodern}
\usepackage{microtype}
\usepackage{amsmath,amssymb,mathtools}
\usepackage{booktabs}
\usepackage[margin=1in]{geometry}
\usepackage{xurl}
\usepackage[hidelinks]{hyperref}
\allowdisplaybreaks[2]
\numberwithin{equation}{section}
\newcommand{\R}{\mathbb R}
\newcommand{\Q}{\mathbb Q}
\newcommand{\sinc}{\operatorname{sinc}}
\newcommand{\diag}{\operatorname{diag}}
\newcommand{\lo}{\mathrm{lo}}
\newcommand{\hi}{\mathrm{hi}}
\newcommand{\eps}{\varepsilon}
\newcommand{\pe}{p_{\mathrm e}}
\newcommand{\pint}{p_{\mathrm i}}
\newcommand{\tS}{\tau_{\mathrm S}}
\newcommand{\tL}{\tau_{\mathrm L}}
\newcommand{\qlo}{\underline q}
\newcommand{\qddlo}{\underline {q''}}
\newcommand{\floor}[1]{\left\lfloor #1\right\rfloor}
\newcommand{\ceil}[1]{\left\lceil #1\right\rceil}
\title[Local certificates and the thirteen-point boundary]{--- Supplementary Information ---\\ Local certificates for compensated Gram bounds:\\
the two seven-point inputs and the completed one-short thirteen-point regions}
\author{R. Arun Chandru}
\date{10 September 2026}
\begin{document}
\begin{abstract}
This supplement gives the mathematical specification
and acceptance proofs for two computer-assisted inequalities on six
nonnegative gaps. They are the cosine row with floor $19/5000$ and
pressure $1/3000$, and the two-mode row with floor $7/400$ and pressure
$1/500$. We explain interval arithmetic, removable singularities,
whole-cell range minima, matrix lower bounds, and complete-domain
accounting, with source-bound records of completed executions.
A separate six-mode, thirteen-point problem is developed from its scalar
exclusions and two analytic baselines. An adjacent-long-gap certificate
and four completed triple bounds then exclude its all-long region and
every one-short region, including both interior placements. The remaining
multi-short region is not certified. No full thirteen-point inequality
or additional global zero-count bound is asserted.
\end{abstract}
\maketitle
\tableofcontents

\section{Scope, dependencies, and the meaning of a certificate}

There are three different Fourier kernels in this supplement. The first
two supply the only local numerical inputs to the main manuscript. The
third concerns a separate, unfinished thirteen-point problem. A floor,
pressure allowance, or quadratic reserve proved for one kernel is not
available for either of the others.

All the mathematics needed to understand the local certificates is given
below. The accompanying verification package is
a computational implementation of these arguments, not a substitute for
them. Its manifest identifies the distributed sources and the provenance
of recorded results. A digest establishes byte identity; by itself it
does not establish an interval enclosure or a covering argument.

The trust boundary has three parts. First are the displayed analytic and
finite-dimensional arguments. Second are a correctly implemented outward
arithmetic backend and the source implementing the stated acceptance and
coverage rules. Third is an actual completed execution of that source
with the exact configuration in question. The recorded executions used
Arb through \texttt{python-flint} version 0.9.0 at 128-bit precision for
scalar cell tables; the later triple certificates use exact rational
arithmetic after these tables have been established. None of these
computations is presented as a formal proof-assistant verification.

The scalar-cell, range-minimum, and convex-tangent verification framework,
including the cosine seven-point row, comes from Ainta's MIT-licensed
implementation \cite{ainta}. The two-mode callbacks and the additional
acceptance rules are adaptations. The Fourier-profile and local-pressure
viewpoint is not claimed here as an invention of this supplement. The
main manuscript discusses its mathematical prior art and the distinct
common-witness transfer. Here the purpose is to expose the precise local
proof obligations and their evidence, including the limits of that
evidence.

\section{The kernels and the local energy}

Put $z=1/\sqrt2$ and define the normalized cosine density by
\begin{equation}\label{eq:density0}
 f_0(t)=\frac{\cos(\sqrt2t)}{\sqrt2\sin z}
 \quad (|t|\leq\tfrac12),\qquad f_0(t)=0\quad(|t|>\tfrac12).
\end{equation}
For a finite real vector $u=(u_1,\ldots,u_d)$, set
\begin{equation}\label{eq:densityu}
 f_u(t)=f_0(t)+\sum_{j=1}^d u_j\cos(2\pi jt)
 \quad (|t|\leq\tfrac12),
\end{equation}
and set it to zero elsewhere. Direct integration gives $\int f_0=1$;
each perturbing cosine has integral zero. Thus $\int f_u=1$.
The three choices used here are
\begin{equation}\label{eq:profiles}
 u_{\mathrm C}=(),\qquad
 u_{\mathrm H}=\left(\frac{13}{500},-\frac7{500}\right),\qquad
 u_{13}=\frac1{4000}(37,-69,-1,8,-11,2).
\end{equation}
The elementary inequalities $\cos z\geq3/4$ and
$0<\sqrt2\sin z\leq1$ imply $f_0\geq3/4$ on its support.
The absolute coefficient sums in the last two profiles are respectively
$1/25$ and $4/125$. Consequently all three densities are nonnegative;
in particular $f_{u_{13}}\geq359/500>0$. Each is real and even.

Our Fourier convention and squared kernel are
\begin{equation}\label{eq:kernel}
 k_u(x)=\int_{-1/2}^{1/2}f_u(t)e^{-2\pi ixt}\,dt,
 \qquad q_u(x)=k_u(x)^2\quad(x\in\R).
\end{equation}
The kernel is real on the real axis, and $q_u\geq0$. Writing
$\sinc y=\sin y/y$ with $\sinc0=1$, direct integration yields the entire
formulas
\begin{align}
 k_0(x)&=\frac{\sinc(\pi x-z)+\sinc(\pi x+z)}{2\sqrt2\sin z},
       \label{eq:sinc0}\\
 k_u(x)&=k_0(x)+\frac12\sum_{j=1}^d u_j
      \{\sinc(\pi(x-j))+\sinc(\pi(x+j))\}.\label{eq:sincu}
\end{align}
These formulas have no actual poles, even when a shifted argument
vanishes. Since the support is compact, differentiation under the
integral is valid to every order. Positivity and normalization give
\begin{equation}\label{eq:globalderivative}
 |k_u|\leq1,\qquad |k_u'|\leq\pi,\qquad |k_u''|\leq\pi^2,
 \qquad |q_u''|\leq4\pi^2.
\end{equation}
The last inequality follows from
$q_u''=2(k_u')^2+2k_uk_u''$.

For $m-1$ nonnegative gaps $g=(g_1,\ldots,g_{m-1})$, define
\begin{equation}\label{eq:Vm}
 V_{m,u}(g)=\sum_{h=1}^{m-1}\frac2{m-h}
       \sum_{i=1}^{m-h}q_u(g_i+\cdots+g_{i+h-1}).
\end{equation}
Zero gaps are included. The coefficients in this formula are part of
the objective and must not be replaced by uniform coefficients over
all pairs. The two seven-point assertions are
\begin{align}
 V_{7,\mathrm C}(g)+\frac1{3000}\sum_{j=1}^6g_j
       &\geq\frac{19}{5000},\label{eq:rowC}\\
 V_{7,\mathrm H}(g)+\frac1{500}\sum_{j=1}^6g_j
       &\geq\frac7{400},\label{eq:rowH}
\end{align}
for every $g\in[0,\infty)^6$. These are whole-orthant statements, not
inequalities only at a mesh or at distinct configurations.

\section{Arithmetic that encloses functions, not samples}

\subsection{Outward operations and removable singularities}

An interval $[a,b]$ denotes a set containing the exact quantity in
question. Sums, differences, products, and quotients are rounded
outwards; division is allowed only after zero has been excluded from
the denominator interval. For rational intervals the endpoints and all
comparisons can be exact. For Arb balls the containment guarantee of the
backend is used, followed, where required, by a downward conversion to
the exact real number represented by a binary64 value. Converting a
lower endpoint to nearest binary64 and then applying
\texttt{nextafter} toward $-\infty$ gives a conservative lower endpoint
when the result is finite. A nonnegative lower bound is clipped at zero
only when nonnegativity has independently been proved.

Away from zero, the identities
\begin{align*}
 \sinc'(y)&=\frac{y\cos y-\sin y}{y^2},\\
 \sinc''(y)&=\frac{(2-y^2)\sin y-2y\cos y}{y^3}
\end{align*}
may be evaluated by interval arithmetic. If an argument ball includes
zero, these divisions are not permitted. The following explicit
polynomials handle such balls. For $|y|\leq b<1$, let
\begin{align*}
 P_0(y)&=\sum_{n=0}^{8}\frac{(-1)^ny^{2n}}{(2n+1)!},\\
 P_1(y)&=-\sum_{n=0}^{8}\frac{(-1)^ny^{2n+1}}{(2n+1)!(2n+3)},\\
 P_2(y)&=-\sum_{n=0}^{8}\frac{(-1)^ny^{2n}}{(2n)!(2n+3)}.
\end{align*}
Then the respective error bounds are
\begin{equation}\label{eq:sincremainders}
 |\sinc-P_0|\leq\frac{b^{18}}{19!},\qquad
 |\sinc'-P_1|\leq\frac{b^{19}}{19!\,21},\qquad
 |\sinc''-P_2|\leq\frac{b^{18}}{18!\,21}.
\end{equation}
For completeness, $\sinc y=\int_0^1\cos(ty)\,dt$.
Differentiate this identity once or twice and apply the real Taylor
remainder for sine or cosine inside the integral. Integrating the
extra powers of $t$ gives exactly the three denominators in
\eqref{eq:sincremainders}. Thus these are bounds over the whole ball,
not an alternating-series test at its center. Substitution in
\eqref{eq:sinc0}--\eqref{eq:sincu}, with factors $\pi$ and $\pi^2$
for differentiation in $x$, encloses $k,k',k''$ and hence $q,q',q''$
for each of the three specified profiles.

\subsection{A completely rational point evaluator}

The scalar gates and the atomic Taylor certificates below can be
checked without floating transcendental arithmetic at their rational
evaluation points. Here is an explicit construction. The Machin
identity
\begin{equation}\label{eq:machin}
 \pi=16\arctan(1/5)-4\arctan(1/239)
\end{equation}
and the sums $\sum_{n=0}^{40}(-1)^n/((2n+1)a^{2n+1})$ give rational
enclosures for the two arctangents. Each error is at most
$1/(83a^{83})$, with its alternating sign. Combine their endpoint
intervals outwards. If desired, after each operation round to the
outward rational lattice of spacing $10^{-50}$.

Put
\begin{equation}\label{eq:R0}
 R_0=\frac12+z\cot z.
\end{equation}
The Taylor sums through index $30$ for $\cos z$ and $\sin z/z$, with
the next alternating term, enclose them using only powers of
$z^2=1/2$. Divide these positive rational intervals to enclose
$z\cot z$. No numerical square root is needed for this operation.

For rational $x$, write $x=n+v$, where $n=\floor{x+1/2}$ and
$-1/2\leq v<1/2$. Evaluate $\sin(\pi v)$ through degree $47$ and
$\cos(\pi v)$ through degree $46$, including symmetric remainders
$2^{49}/49!$ and $2^{48}/48!$. These follow from $|\pi v|<2$ and
the ordinary real Taylor theorem. Interval evaluation of these
polynomials, and the factor $(-1)^n$, encloses $\sin(\pi x)$ and
$\cos(\pi x)$.

The directly integrated point formulas are
\begin{align}
 k_0(x)&=\frac{\cos(\pi x)-2(R_0-1/2)\pi x\sin(\pi x)}
                    {1-2\pi^2x^2},\label{eq:rationalkernel}\\
 k_u(x)-k_0(x)&=\sum_{j=1}^d
       u_j\frac{(-1)^jx\sin(\pi x)}{\pi(x^2-j^2)}.
       \label{eq:rationalmodes}
\end{align}
At $x=j$ the corresponding mode integral is exactly $1/2$ and is
handled separately. Other denominators must be checked to exclude zero.
All rational point evaluations below have $x\geq5/6$. The atomic
helper includes the separately guarded range $[9/10,1)$.
On that below-one range $1-2\pi^2x^2<0$ and each $x^2-j^2<0$; the interval
implementation still verifies these signs before division. At the
positive integers the removable resonances are handled explicitly.
This is a point evaluator with proved error bounds. A whole interval is
enclosed either by the entire sinc formulas or by an additional
derivative bound; point evaluation alone is never a domain certificate.

\section{The complete seven-point verification}

\subsection{Pressure compactification and scalar cell tables}

Write $F_7=V_7+p\sum g_j$, with one of the exact pairs
$(\eps,p)$ in \eqref{eq:rowC}--\eqref{eq:rowH}. Since $V_7\geq0$,
the region $\sum g_j\geq\eps/p$ is settled by pressure alone.
Every strict counterexample would therefore have all its coordinates
less than $\eps/p$. The common grid is $G=4000$, with closed cells
\begin{equation}\label{eq:cells}
 J_i=[i/G,(i+1)/G]\qquad(i=0,1,\ldots).
\end{equation}
Let $L=\ceil{\eps G/p}$. The coordinate cells $0,\ldots,L-1$
cover every possible strict counterexample; shared boundaries are
included. The recorded tables have $L+8$ entries. They supply lower
bounds
\begin{equation}\label{eq:tables}
 q(x)\geq Q_i\geq0,\qquad q''(x)\geq D_i\quad(x\in J_i)
\end{equation}
where the second assertion is made only for finite, available entries.
To obtain $Q_i$, enclose $k(J_i)$ by Arb, take a nonnegative lower
bound for its absolute value, convert downwards, and square downwards.
To obtain $D_i$, enclose $q''(J_i)$ and convert its lower endpoint
downwards. An unavailable curvature prefix is represented by
$-\infty$, not zero.

It is unnecessary for these tables to contain every possible summed
distance in the six-dimensional box. If a sum reaches beyond the
stored $q$ table, the universally valid lower bound zero is used for
that energy term. If it reaches beyond the stored curvature table,
the corresponding derivative test is unavailable. No extrapolation or
unproved zero-curvature assertion is made.

Before forming multidimensional roots, remove any coordinate cell for
which
\begin{equation}\label{eq:singleprune}
 \frac13 Q_i+p\frac iG\geq\eps.
\end{equation}
This is valid because that coordinate's nearest-gap term in $V_7$
has coefficient $1/3$, and every discarded term of $F_7$ is
nonnegative. Maximal consecutive surviving groups form scalar
components. Their Cartesian products cover all remaining candidates.

\subsection{Range indices and elementary lower bounds}

Represent a box by six inclusive integer ranges $[a_j,b_j]$.
Its physical coordinate interval is
$[a_j/G,(b_j+1)/G]$. A sum of $h$ consecutive coordinates is
contained in the scalar cells from
\begin{equation}\label{eq:sumindices}
 a_i+\cdots+a_{i+h-1}
 \quad\hbox{through}\quad
 b_i+\cdots+b_{i+h-1}+h-1,
\end{equation}
inclusive. The final $h-1$ is essential. Omitting it would leave
part of the closed sum interval unchecked.

Taking the minimum of every $Q$ entry in this range gives a whole-box
lower bound for its energy term. A sparse range-minimum structure can
answer the same query using two overlapping power-of-two ranges whose
union is the entire requested range. Overlap is harmless for a
minimum; a gap would not be. Adding all the resulting lower bounds
with the original positive coefficients and the pressure at lower
coordinate endpoints gives a valid interval lower bound for $F_7$.
The target is rounded upwards and lower contributions downwards
before any binary64 acceptance comparison.

\subsection{Signed Hessian bounds and center tests}

For each pair distance write $x_e=v_e^Tg$, where $v_e$ is the zero-one
incidence vector of its consecutive gaps, and let $c_e=2/(7-h_e)$,
where $h_e=v_e^Tv_e$. The exact Hessian is
\begin{equation}\label{eq:hessian7}
 \nabla^2 F_7(g)=\sum_e c_e q''(v_e^Tg)v_ev_e^T.
\end{equation}
If $\ell_e\leq c_eq''(v_e^Tg)$ throughout a box, then the exact
represented matrix
\begin{equation}\label{eq:matrixlower}
 H_-:=\sum_e\ell_ev_ev_e^T\preceq\nabla^2F_7(g)
\end{equation}
throughout it. Indeed, the difference is a sum of positive
semidefinite rank-one matrices with nonnegative coefficients. Negative
$\ell_e$ are allowed. When $c_e$ is bracketed by binary64 endpoints,
its lower endpoint must be used for a nonnegative curvature lower
bound and its upper endpoint for a negative one; the product is then
rounded downwards. This sign rule prevents an upward error in a
negative Hessian contribution.

Let $c$ be the rational center of the physical box, and let $r_j$
be its rational coordinate radii. If rigorous LDL decomposition
establishes $H_-$ positive definite, $F_7$ is convex on the box.
In particular,
\begin{equation}\label{eq:tangent}
 F_7(g)\geq F_7(c)-\sum_{j=1}^6|\partial_j F_7(c)|r_j.
\end{equation}
The center value is bounded below and each derivative magnitude above
by Arb. A floating LDL test may be used to decide whether to attempt
the rigorous test, but it cannot accept a box. The rigorous LDL test
uses the exact represented rational matrix and requires every pivot
strictly positive, with all interval divisions justified.

The two-mode computation also uses the stronger lower quadratic.
Taylor's integral remainder along the segment from $c$ to $g$ gives
\begin{equation}\label{eq:quadratic7}
 F_7(g)\geq F_7(c)+\nabla F_7(c)^T(g-c)
                   +\tfrac12(g-c)^TH_-(g-c).
\end{equation}
If $H_-=LDL^T$ is positive definite, minimizing the right-hand side
over all of $\R^6$ gives
\begin{equation}\label{eq:debit7}
 F_7(g)\geq F_7(c)-\tfrac12\nabla F_7(c)^TH_-^{-1}\nabla F_7(c).
\end{equation}
Solving $Ly=\nabla F_7(c)$ by interval forward substitution encloses
the quadratic term as $\sum_j y_j^2/D_j$. It is subtracted with an
upper bound. The minimizing point need not belong to the box.

A further valid fallback discards positive curvature contributions.
For $\ell_e<0$, Cauchy--Schwarz gives
$\ell_ev_ev_e^T\succeq\ell_e h_e I$.
Thus $\mu I\preceq H_-$ with
$\mu=\sum_{\ell_e<0}h_e\ell_e\leq0$. The lower Taylor bound becomes
\begin{equation}\label{eq:signedfallback}
 F_7(g)\geq F_7(c)-\sum_j|\partial_jF_7(c)|r_j
                       +\frac\mu2\sum_jr_j^2.
\end{equation}
All quantities are enclosed outwards. Each of
\eqref{eq:tangent}, \eqref{eq:debit7}, and
\eqref{eq:signedfallback} is an alternative lower bound. Their
maximum is permissible; adding them would count the same cost more
than once.

Center values may be cached because every center distance is exactly
of the form $k/(2G)$ for an integer $k$. A cache key must encode this
exact index, and the profile and precision must be immutable during
the run. Caching complete enclosing balls does not change any
acceptance rule.

\subsection{Exhaustion of the domain}

If a box passes none of the valid lower bounds, bisect its widest
integer coordinate range. For $a<b$, the two children $[a,m]$ and
$[m+1,b]$, $m=\floor{(a+b)/2}$, partition its atomic cell indices
and cover its closed physical interval. Continue until all roots have
been discharged. A remaining atomic box that passes no test is
unresolved, not accepted and not automatically a counterexample.
Similarly, a node cap, exception, or interrupted execution is not a
successful proof.

Uniform pressure and \eqref{eq:Vm} are invariant under full reversal
of the six gaps. A reversal reduction, if used, must keep one
representative for each genuinely reversed pair. There is no
arbitrary permutation symmetry. Source review and complete tree
traversal establish coverage; identities such as
\begin{equation}\label{eq:treecounts}
 \text{roots}+\text{splits}=\text{accepted leaves},\qquad
 \text{visited nodes}=\text{splits}+\text{accepted leaves}
\end{equation}
are useful independent consistency checks on a completed binary forest,
but are not a replacement for that coverage argument.

\section{Completed verification of the two seven-point rows}

The distributed verification package recomputed both seven-point
inequalities with the exact configurations above. The following table
records these completed executions. Both use grid 4000, 128-bit Arb
arithmetic, and curvature start index 1001. Each successful execution
exhausted its covering forest with no unresolved atomic cell.

\begin{center}\small
\begin{tabular}{@{}lrr@{}}
\toprule
Quantity & Cosine row & Two-mode row\\
\midrule
Exact floor & $19/5000$ & $7/400$\\
Pressure per gap & $1/3000$ & $1/500$\\
Grid / Arb precision & $4000/128$ & $4000/128$\\
Coordinate cutoff $L$ & $45600$ & $35000$\\
Scalar table entries & $45608$ & $35008$\\
Curvature start index & $1001$ & $1001$\\
Initial roots & $378$ & $1$\\
Visited nodes & $268856$ & $129093$\\
Splits & $134239$ & $64546$\\
Accepted leaves & $134617$ & $64547$\\
Pressure leaves & $1597$ & $602$\\
Interval leaves & $111698$ & $38533$\\
Derivative-test leaves & $21322$ & $25412$\\
Maximum depth & $32$ & $59$\\
Unresolved / pending & $0/0$ & $0/0$\\
\bottomrule
\end{tabular}
\end{center}

The cosine coordinate components, in inclusive cell indices, are
\begin{equation}\label{eq:Ccomponents}
 [3809,4778],\qquad[7221,9363],\qquad[10572,44827].
\end{equation}
There are $3^6=729$ Cartesian roots, of which
$(729+27)/2=378$ represent the reversal orbits. The
two-mode computation has the single component $[3443,34516]$ and
therefore one six-dimensional root. Its category named ``tangent''
in the saved report includes the valid derivative alternatives above;
that label is not a claim that each such leaf used only the affine
tangent formula. Its center-cache report has $1666485$ hits and $2574$
misses.

The scalar-table SHA-256 values, with binary64 table encoding fixed
by the accompanying source and manifest, are as follows. Line breaking
of a printed digest has no mathematical significance.
\begin{quote}\footnotesize
Cosine $q$:
\path{a9992300d2bf71665aa2b6bd2727e798624cd297103bb200c7f0ca2baea55a2c}

Cosine $q''$:
\path{766b97d9ba81dffd915acaf1c1ba9938166c3191011e744424c12fd3a141bd52}

Two-mode $q$:
\path{f0be0477f6c4c98c7c3d879abf79f509fa637e9d4740fe8ee83bd4af52899580}

Two-mode $q''$:
\path{aa35a51991d4a84c748e1908f7d95482ac507051d7c67b16c6919ee09f35377c}
\end{quote}

The cosine row has a completed ordinary execution of the distributed
source. The two-mode row has completed ordinary and \texttt{-OO}
executions: their entire receipts agree after removing only the
optimization flag and the engine's elapsed time. Source identities,
table digests and machine-readable covering summaries accompany the
package. These completed computations support \eqref{eq:rowC} and
\eqref{eq:rowH} subject to the stated implementation and arithmetic
trust boundary. An optimization-mode comparison is a consistency
check, not an independent enclosure method.

\section{A separate thirteen-point objective}

For the remainder of the supplement, $k=k_{u_{13}}$ and $q=k^2$
refer only to the six-mode profile in \eqref{eq:profiles}. Define
\begin{equation}\label{eq:F13}
 F(g)=V_{13,u_{13}}(g)+\sum_{j=1}^{12}p_jg_j,
 \qquad g\in[0,\infty)^{12},
\end{equation}
where
\begin{equation}\label{eq:pressures13}
 p_1=p_{12}=\pe=\frac{811}{4800000},\qquad
 p_2=\cdots=p_{11}=\pint=\frac{6359}{24000000}.
\end{equation}
In particular,
\begin{equation}\label{eq:target13}
 P:=\sum_jp_j=\frac{239}{80000},\qquad
 \eps:=\frac{31}{5000}.
\end{equation}
We do \emph{not} assert $F\geq\eps$ on this entire orthant. The
first task is to restrict where a strict counterexample could lie.
All scalar restrictions in the next two sections are necessary
conditions under the explicit hypothesis $F(g)<\eps$.

Write $a=1/(\sqrt2\sin z)$. Since
$11/12\leq\sin z/z\leq1$, we have $1\leq a\leq12/11$.
The convenient full-density bounds
\begin{equation}\label{eq:densitybounds13}
 m:=\frac{71}{100}\leq f_{u_{13}}(t)\leq
 M:=\frac{311}{275}\quad(|t|\leq\tfrac12)
\end{equation}
follow from $\cos z\geq3/4$ and
$\sum|u_j|=4/125\leq1/25$. The bounds apply to the perturbed
density, not just its cosine part.

\subsection{Monotonicity and the minimum-gap bootstrap}

Differentiating the even Fourier integral gives
\begin{equation}\label{eq:kmonotone}
 -k'(x)=4\pi\int_0^{1/2}t f_{u_{13}}(t)\sin(2\pi xt)\,dt.
\end{equation}
For $0<x\leq1$ this is strictly positive. For
$1\leq x\leq27/25$, the sine is at least $1/2$ on
$[1/8,3/8]$: its argument belongs to
$[\pi/4,81\pi/100]\subset[\pi/6,5\pi/6]$.
This part of the integral in \eqref{eq:kmonotone}, before its
factor $4\pi$, is at least $71/3200$. A negative contribution can
occur only for $t\geq1/(2x)\geq25/54$; its absolute value is at
most $M(1/2)(1/2-25/54)=311/14850$. Their difference is
$1183/950400>0$. All other contributions are nonnegative. Hence
$k$ is strictly decreasing on $(0,27/25]$.

The exact point evaluator in Section 3 verifies
$k(5/6)>0$ and $k(5/6)^2/6>\eps$.
If some gap were at most $5/6$, monotonicity would make its nearest
term alone exceed the target. Thus all gaps of a strict
counterexample exceed $5/6$. More generally, suppose all such gaps
already exceed $d_0$, and $d_0<d_1\leq1$ with $k(d_1)>0$.
A gap at most $d_1$ would imply
\begin{equation}\label{eq:bootstrap1}
 F(g)\geq\frac{k(d_1)^2}{6}+d_0P.
\end{equation}
The following outward lower endpoints exceed $\eps$ and prove the
successive improved cuts. The pressure $d_0P$ includes all twelve
gaps exactly once; only one nonnegative nearest energy term is kept.
\begin{center}\small
\begin{tabular}{@{}lll@{}}
\toprule
$d_0$ & $d_1$ & Lower endpoint in \eqref{eq:bootstrap1}\\
\midrule
$5/6$ & $9/10$ & $0.007009907197317624826482$\\
$9/10$ & $91/100$ & $0.006608495667780589234137$\\
$91/100$ & $183/200$ & $0.006356600267380480502615$\\
$183/200$ & $917/1000$ & $0.006262235820193320264541$\\
$917/1000$ & $459/500$ & $0.006214284540455943193710$\\
\bottomrule
\end{tabular}
\end{center}

Set $\tau_0=209/200$. Exact evaluation encloses
\begin{equation}\label{eq:tau0value}
 k(\tau_0)\in
 [0.015366680647566453432485,\ 0.015366680647566453432486]
\end{equation}
and gives
\begin{equation}\label{eq:tau0reserve}
 \frac{k(\tau_0)^2}{6}-\pint(\tau_0-459/500)
       >0.000005706104020715559434.
\end{equation}
For $459/500\leq g\leq\tau_0$, monotonicity, positivity, and
\eqref{eq:tau0reserve} imply
$q(g)/6+pg\geq p\tau_0$ for either pressure $p$.
For $g\geq\tau_0$ pressure alone proves the same assertion.
Thus this first scalar baseline holds on the entire half-line
$[459/500,\infty)$.

If all gaps exceed $d_0$ but one is at most $d_1\leq\tau_0$,
use the eleven other scalar baselines to obtain
\begin{equation}\label{eq:bootstrap2}
 F(g)\geq\frac{k(d_1)^2}{6}+\tau_0P-\pint(\tau_0-d_0).
\end{equation}
The steps $(d_0,d_1)=(.918,.920),(.920,.922),(.922,.924)$
have respective outward lower endpoints
\begin{align*}
 &0.006456625078131788200580,\\
 &0.006352660139980702315412,\\
 &0.006250599414089590385868.
\end{align*}
They exceed $\eps$. Consequently every strict counterexample has
\begin{equation}\label{eq:delta}
 g_j>\delta:=\frac{231}{250}\quad(1\leq j\leq12).
\end{equation}
Each decimal just displayed is a terminating rational lower bound,
not a claim about an unbounded numerical error.

\subsection{Compact necessary domains}

Put
\begin{equation}\label{eq:gateB}
 B=\eps-\tau_0P=\frac{49249}{16000000},\qquad
 W_p(x)=\frac{q(x)}6+p(x-\tau_0).
\end{equation}
The nearest cost $w_p(x)=q(x)/6+px$ is at least both $p\tau_0$
and $px$. Since all higher-distance terms are nonnegative, a strict
counterexample satisfies
\begin{equation}\label{eq:positiveparts0}
 \sum_jp_j(g_j-\tau_0)_+<B,
 \qquad W_{p_j}(g_j)<B.
\end{equation}
The second inequality follows by using the other eleven baselines.
In particular its coordinates and span satisfy
\begin{align}
 g_1,g_{12}&<\tau_0+B/\pe=\frac{3124439}{162200},\nonumber\\
 g_2,\ldots,g_{11}&<\tau_0+B/\pint=\frac{16103731}{1271800},
       \label{eq:globalcaps}\\
 \sum_jg_j&<12\tau_0+B/\pe=\frac{623616}{20275}.
       \nonumber
\end{align}
The final assertion uses
$\sum g_j\leq12\tau_0+\sum(g_j-\tau_0)_+$ and
$p_j\geq\pe$. It does not assume the gaps lie close to one or two.

\section{The exact one-dimensional gate}

\subsection{A point enclosure plus an analytic Lipschitz bound}

For $x>0$, integration by parts in the Fourier integral yields
\begin{equation}\label{eq:ibp}
 k(x)=\frac{f(1/2)\sin(\pi x)-
              \int_0^{1/2}f'(t)\sin(2\pi xt)\,dt}{\pi x}.
\end{equation}
Here $f=f_{u_{13}}$ on its support; the endpoint contribution is
explicit, so no smooth extension by zero is assumed. The cosine
part decreases on $[0,1/2]$. Hence
\begin{align*}
 |f(1/2)|+\int_0^{1/2}|f'(t)|\,dt
 &\leq a+\sum_j|u_j|+2\sum_jj|u_j|\\
 &\leq C:=\frac{27751}{22000}.
\end{align*}
For the perturbation this uses
$\int_0^{1/2}|2\pi j\sin(2\pi jt)|\,dt=2j$.
It follows that $|k(x)|\leq C/(\pi x)$. Differentiating
\eqref{eq:ibp} gives
\[
 k'(x)=\frac{f(1/2)\cos(\pi x)}x-\frac{k(x)}x
              -\frac2x\int_0^{1/2}t f'(t)\cos(2\pi xt)\,dt,
\]
so
\begin{equation}\label{eq:decaybounds}
 |k'(x)|\leq\frac Cx+\frac C{\pi x^2}.
\end{equation}
The integral formula and $f\leq M$ also give
$|k'(x)|\leq\pi M/2$.

For a closed rational interval $I=[l,r]$ with $l>0$, let
$c=(l+r)/2$, $h=(r-l)/2$, and let
$[k_{\lo},k_{\hi}]$ enclose $k(c)$ by Section 3. For an outward
interval $[\pi_{\lo},\pi_{\hi}]$, define
\begin{equation}\label{eq:KL}
 K=\min\left(1,\frac C{\pi_{\lo}l}\right),\qquad
 L=\min\left(\frac{\pi_{\hi}M}{2},
                   \frac Cl+\frac C{\pi_{\lo}l^2}\right).
\end{equation}
Then $k(I)\subset[v,w]=[k_{\lo}-Lh,k_{\hi}+Lh]$.
Put
\begin{equation}\label{eq:absolutegate}
 A_{\lo}=\max(0,v,-w),\qquad
 A_{\hi}=\min(K,\max(|v|,|w|)).
\end{equation}
These bound $|k|$ below and above on $I$. Therefore
\begin{equation}\label{eq:Wgate}
 \frac{A_{\lo}^2}{6}+p(l-\tau_0)
 \ \leq\ W_p(x)\ \leq\
 \frac{A_{\hi}^2}{6}+p(r-\tau_0)\qquad(x\in I).
\end{equation}
Every operation here is on rational endpoints with outward rounding.
This is a whole-interval argument derived from the derivative bound,
not midpoint sampling.

\subsection{Closed partition and retained components}

For each of the two pressures, start with the closed interval
$[\delta,\tau_0+B/p]$. If the lower bound in
\eqref{eq:Wgate} is at least $B$, exclude the whole interval.
Equality is sufficient because the sublevel under consideration is
strict. If its upper bound is less than $B$, retain the interval.
Otherwise bisect it unless its width is at most $1/16384$; such
unresolved boundary intervals are retained. Any emergency budget
exhaustion must also retain pending intervals. Neither uncertainty
nor termination may silently discard a possible counterexample.

The distributed gate checker completed both closed partitions.
The endpoint gate has 50 excluded, 72 strictly passing,
and 7 unresolved terminal intervals. The interior gate has 55,
74, and 10 respectively. No emergency budget was exhausted.
Across both pressures, each run visited 534 cells and used 519
distinct exact midpoint values. A preliminary envelope-only test,
using $0\leq q\leq K^2$, can settle a cell without a point evaluation.
Exact sorting, shared
endpoints, and total interval lengths checked that the final leaves
partitioned each initial domain. The ordinary and \texttt{-OO}
executions agree in their mathematical output.

Merging retained adjacent leaves gives two closed necessary components
for each pressure. Their exact endpoints are recorded here so that
no floating approximation is used to define a domain:
\begin{align}
 S_{\mathrm e}^{\rm gate}
 &=\left[\frac{393403101917}{425197568000},
         \frac{538442949829}{425197568000}\right],\nonumber\\
 L_{\mathrm e}^{\rm gate}
 &=\left[\frac{183665052303}{106299392000},
         \frac{8164814247361}{425197568000}\right],\label{eq:gates}\\
 S_{\mathrm i}^{\rm gate}
 &=\left[\frac{770962919881}{833486848000},
         \frac{1054606088081}{833486848000}\right],\nonumber\\
 L_{\mathrm i}^{\rm gate}
 &=\left[\frac{22558120767}{13023232000},
         \frac{20694035057199}{1666973696000}\right].\nonumber
\end{align}
The following slightly enlarged domains suffice for the local
arguments:
\begin{equation}\label{eq:enlargeddomains}
 S_{\mathrm e}=[459/500,317/250],\qquad
 S_{\mathrm i}=[459/500,633/500],\qquad L=[5/3,\infty).
\end{equation}
The inclusions of \eqref{eq:gates} in the corresponding sets of
\eqref{eq:enlargeddomains} are exact rational comparisons. Thus a
strict counterexample has an S or L label at every coordinate.
This does not label every arbitrary nonnegative twelve-tuple: the
strict-counterexample hypothesis is indispensable.

\section{Analytic short and long baselines, with reserves}

\subsection{The short baseline $53/50$}

We first prove, independently of any gate, that
\begin{equation}\label{eq:shortbaseline}
 \frac{q(x)}6+px\geq\tS p,\qquad \tS=\frac{53}{50},
 \qquad x\geq459/500,
\end{equation}
for either pressure $p$.
For $459/500\leq x\leq1$, the sine in
\eqref{eq:kmonotone} is nonnegative. Exact integration gives
\begin{align*}
 -k'(x)&\geq m\frac{\sin(\pi x)-\pi x\cos(\pi x)}{\pi x^2}\\
 &\geq m\frac{-\cos(\pi x)}x
 \geq m\cos(\pi/10)>\frac{19m}{20}>\frac35.
\end{align*}
Here $x\geq.9$, $x\leq1$, and
$\cos(\pi/10)\geq1-\pi^2/200>19/20$ using $\pi^2<10$.
For $1\leq x\leq T:=1063/1000$, split the integral at
$b=1/(2x)$. On $[0,b]$ the positive part is at least
$m/(4\pi x^2)$. On $[b,1/2]$, use
$|\sin(2\pi xt)|\leq2\pi x(t-b)$ and $t\leq1/2$ to bound
the absolute negative part above by $M\pi(x-1)^2/(8x)$.
Multiplying by $4\pi$ gives
\begin{align}
 -k'(x)&\geq\frac m{x^2}-\frac{M\pi^2(x-1)^2}{2x}
 \geq\frac m{T^2}-5M(T-1)^2\nonumber\\
 &=\frac35+\frac{366235595129}{62148295000000}>\frac35.
 \label{eq:shortderivative}
\end{align}
This proves $-k'>3/5$ throughout $[459/500,1063/1000]$.

The exact point evaluation
\begin{equation}\label{eq:shortpoint}
 k(53/50)\in
 [0.001990173438104262028271,\ 0.001990173438104262028272]
\end{equation}
has lower endpoint greater than both $a_S:=19/10000$ and
$5\pint$. On $[459/500,\tS]$ we therefore have
$k(x)\geq k(\tS)>0$ and
\[
 \frac{d}{dx}\{q(x)/6+px\}=k(x)k'(x)/3+p
       \leq-k(\tS)/5+p<0.
\]
The nearest cost is at least its value at $\tS$, which is at
least $p\tS$. For $x\geq\tS$, pressure alone proves
\eqref{eq:shortbaseline}. No derivative assertion beyond the
stated finite interval is used.

There is also a useful explicit reserve. For
$459/500\leq s\leq\tS$ put $d=\tS-s$. Integrating
\eqref{eq:shortderivative} gives $k(s)\geq a_S+(3/5)d$.
Consequently
\begin{equation}\label{eq:shortreserve}
 r_{S,p}(s):=q(s)/6+p(s-\tS)
 \geq\frac{a_S^2}{6}+(a_S/5-p)d+\frac3{50}d^2.
\end{equation}
All three coefficients are positive for both pressures.

\subsection{The long baseline $203/100$}

Put $A=\sqrt2\cot z=2z\cot z$. The same elementary sine and
cosine estimates give $3/2\leq A\leq24/11$. Differentiating
\eqref{eq:rationalkernel}, with $v=\pi x$, yields
\begin{equation}\label{eq:k0longderivative}
 k_0'(x)=\frac{\pi}{(2v^2-1)^2}
 \left[(2Av^3+(4-A)v)\cos v
       -(2(A-1)v^2+1+A)\sin v\right].
\end{equation}
On $5/3\leq x\leq41/20$ we have $v>5$, $\cos v\geq1/2$,
$\sin v\leq\pi/20$, and $4-A>0$. The bracket is therefore at
least
$Q(v)=Av^3-(\pi/20)(2(A-1)v^2+1+A)$.
Using $\pi^2<10$, $A-1\leq13/11$, $1+A\leq35/11$, and $v>5$
gives
\begin{equation}\label{eq:Qbound}
 \frac{\pi Q(v)}{4v^4}
 \geq\frac{15}{82}-\frac12\left(\frac{13}{550}
                                      +\frac{35}{27500}\right)
 =\frac{15}{82}-\frac{137}{11000}>0.
\end{equation}
For the first term we used
$\pi A/(4v)=A/(4x)\geq15/82$. Since $Q>0$, replacing
$(2v^2-1)^2$ by the larger $4v^4$ in the positive lower bound is
legitimate. Hence \eqref{eq:Qbound} is a lower bound for $k_0'$.

The perturbation has absolute coefficient sum $4/125$, and therefore
\[
 |(k-k_0)'(x)|\leq2\pi\int_{-1/2}^{1/2}|t|
       \sum_j|u_j|\,dt=\frac{2\pi}{125}<\frac{44}{875}.
\]
Combining this with \eqref{eq:Qbound} proves
\begin{equation}\label{eq:longderivative}
 k'(x)\geq\frac{15}{82}-\frac{137}{11000}-\frac{44}{875}
       =\frac3{25}+\frac{589}{3157000}>\frac3{25}
 \quad(5/3\leq x\leq41/20).
\end{equation}
The exact point evaluation is
\begin{equation}\label{eq:longpoint}
 k(203/100)\in
 [-0.008907305837779524800319,\ -0.008907305837779524800318],
\end{equation}
so $k(\tL)<-1/125$ for $\tL:=203/100$.
For $5/3\leq x\leq\tL$, put $d=\tL-x$ and integrate
\eqref{eq:longderivative} to get
$k(x)\leq-1/125-(3/25)d$. Thus
\begin{equation}\label{eq:longreserve}
 r_{L,p}(x):=q(x)/6+p(x-\tL)
 \geq\frac1{93750}+(1/3125-p)d+\frac3{1250}d^2.
\end{equation}
The linear coefficient is positive, since
$1/3125-\pint=1321/24000000>0$. For $x\geq\tL$,
nonnegativity of $q$ and pressure suffice. We have proved the entire
half-line baseline
\begin{equation}\label{eq:longbaseline}
 q(x)/6+px\geq\tL p\qquad(x\geq5/3)
\end{equation}
for both pressures. Neither this proof nor the short proof uses
derivative samples.

\subsection{Pattern budgets and safe use of reserves}

Fix any S/L label pattern permitted by the gates, and put $t_j=\tS$
on S coordinates and $t_j=\tL$ on L coordinates. Define
\begin{equation}\label{eq:patternAD}
 A_{\rm pat}=\sum_jp_jt_j,\qquad
 \Delta_{\rm pat}=\eps-A_{\rm pat}.
\end{equation}
For a strict counterexample, the two baselines and $w_{p_j}(g_j)\geq
p_jg_j$ imply
\begin{equation}\label{eq:patterncaps}
 \sum_jp_j(g_j-t_j)_+<\Delta_{\rm pat},\qquad
 g_j<t_j+\frac{\Delta_{\rm pat}}{p_j},\qquad
 \sum_jg_j<\sum_jt_j+\frac{\Delta_{\rm pat}}{\pe}.
\end{equation}
If there are $a$ long interior and $b$ long endpoint coordinates,
these constants have the closed form
\begin{align}
 A_{a,b}&=\frac{53}{50}P+\frac{97}{100}(a\pint+b\pe),
       \nonumber\\
 \Delta_{a,b}&=\frac{12133}{4000000}
                   -\frac{97}{100}(a\pint+b\pe),\label{eq:countbudgets}\\
 \sum_jt_j&=12\frac{53}{50}+\frac{97}{100}(a+b).\nonumber
\end{align}
All 33 budgets, $0\leq a\leq10$, $0\leq b\leq2$, are positive:
the largest baseline is the all-long baseline, whose deficit is
$1083/8000000>0$. Thus these scalar budgets alone exclude no complete
label pattern, and they do not make arrangements with the same
counts equivalent for the longer-distance energy terms.

There is a useful optional closed span bound from the upper gate
endpoints $U_{eS},U_{eL},U_{iS},U_{iL}$ in \eqref{eq:gates}.
Let
\begin{align*}
 Z_e&=(2-b)(U_{eS}-\tS)+b(U_{eL}-\tL),\\
 Z_i&=(10-a)(U_{iS}-\tS)+a(U_{iL}-\tL),\\
 z_e&=\min(Z_e,\Delta/\pe),\qquad
 z_i=\min(Z_i,\max(\Delta-\pe z_e,0)/\pint).
\end{align*}
Then $\sum g_j\leq\sum t_j+z_e+z_i$. Indeed, relax the
positive excesses $v_j=(g_j-t_j)_+$ to
$0\leq v_j\leq U_j-t_j$, $\sum p_jv_j\leq\Delta$.
Since $\pe<\pint$, exchanging interior excess for available
endpoint excess improves the span at no greater cost. Spend the
budget on the two endpoint coordinates first, then the interiors;
this gives exactly $z_e+z_i$. This bound is stated non-strictly.
It does not authorize an equality rejection appropriate only to the
strict caps in \eqref{eq:patterncaps}.

The polynomial reserves \eqref{eq:shortreserve} and
\eqref{eq:longreserve} bound the nearest terms including their
pressures. On a box they may replace a weaker lower bound for that
same nearest term. They must not be added to an energy estimate that
has already spent it. If all coordinates except $j$ have nonnegative
residual reserves $\sigma_\ell$, a valid further necessary cap is
$g_j<t_j+(\Delta-\sum_{\ell\ne j}\sigma_\ell)/p_j$.
The reserve for coordinate $j$ is deliberately absent from the sum.

\section{One adjacent-long-gap credit}

\subsection{A pressure choice dominated by both actual choices}

For $x\geq5/3$ define
\begin{equation}\label{eq:worstlong}
 \phi_L(x)=
 \begin{cases}
 \pint(x-\tL),&x\leq\tL,\\
 \pe(x-\tL),&x\geq\tL,
 \end{cases}
 \qquad r_L^*(x)=q(x)/6+\phi_L(x).
\end{equation}
This is the minimum of the two actual long residuals. Hence it is
nonnegative by \eqref{eq:longbaseline} and no greater than either
$r_{L,p}(x)$. The function $\phi_L$ is increasing and continuous,
with its only slope change at $\tL$.

The adjacent-long credit is the whole-domain inequality
\begin{equation}\label{eq:LLcredit}
 r_L^*(x)+r_L^*(y)+\frac2{11}q(x+y)
       \geq K_L:=\frac3{100000}\qquad(x,y\geq5/3).
\end{equation}
Here is its analytic compactification and finite certificate.
Let
\begin{equation}\label{eq:longP}
 P_L(d)=\frac1{93750}+\frac{1321}{24000000}d+\frac3{1250}d^2.
\end{equation}
It increases for $d\geq0$. If $x\leq39/20$, then
$d=\tL-x\geq2/25$ and
$r_L^*(x)\geq P_L(2/25)=3043/10^8>K_L$.
If $x\geq\tL+K_L/\pe=179033/81100$, its pressure residual
alone pays $K_L$. The same applies to $y$. It remains to cover
the closed square
\begin{equation}\label{eq:LLcore}
 [39/20,179033/81100]^2.
\end{equation}

\subsection{The complete scalar tables and finite square}

For the six-mode profile the distributed generator completed scalar
tables using grid $G=500$,
128-bit precision, and 15393 entries. The unused prefix below index
462 has nonnegative lower bound zero and unavailable curvature
$-\infty$; entries from 462 onward were computed on the entire
closed cells using Section 3. The complete-table digests are
\begin{quote}\footnotesize
$q$:
\path{720514E4C09259C32629BBA740B5F7E3D61A3E85F6389AD7AE27BC61C0BB305C}

$q''$:
\path{A67323D95F9407B91ACF10BB517C60D8F4B2BB9CB73DA8E4AFC2523630EA0CB2}.
\end{quote}
The package includes the generated tables and their source-bound
generation record. Its format specifies the exact profile, closed-cell
convention, entry count, precision, unavailable prefix and byte order.
The generator rebuilds the enclosures; matching a stored digest alone
does not establish them. No twelve-dimensional covering claim is made
by this scalar calculation.

The square \eqref{eq:LLcore} is covered outwards by cell indices
$975,\ldots,1103$, giving 129 choices for each coordinate. Let
$Q_i$ be its exact represented binary64 lower bound for $q$ and put
$R_i=\max(0,Q_i/6+\phi_L(i/500))$. The maximum with zero is valid
because the true residual is nonnegative. For cells $(i,j)$ the
test is the exact rational comparison
\begin{equation}\label{eq:LLcell}
 R_i+R_j+\frac2{11}\min(Q_{i+j},Q_{i+j+1})\geq K_L.
\end{equation}
Both sum cells are present; they cover the full interval for $x+y$.
The expression is symmetric, so all $129\cdot130/2=8385$
unordered pairs suffice. Each of the four triple invocations below
recomputed all 8385 comparisons before using an LL replacement, both
in ordinary and optimized Python. The smallest lower bound
was
\begin{equation}\label{eq:LLminimum}
 \frac{23653871141765599661283}{634106827533765836800000000}
       >\frac3{100000}.
\end{equation}
Together with the closed tails, this proves \eqref{eq:LLcredit}
subject to the completed-table arithmetic boundary described above.

A use of \eqref{eq:LLcredit} replaces its two nearest residuals and
their connecting lag-two energy term by $K_L$. It does not add
$K_L$ while retaining those terms unchanged. In a longer string,
such replacements can be made on disjoint adjacent pairs. For
example, maximum matching credits $w_j$ along a path can be
selected by the elementary recurrence
$D_0=D_1=0$, $D_v=\max(D_{v-1},D_{v-2}+w_{v-1})$.
Two neighboring LL replacements sharing a middle residual may not
be added. Only this disjoint replacement principle is needed below.

\section{Four triple inequalities and their finite cores}

For labels $T_1,T_2,T_3\in\{L,S\}$ define the local triple cost
\begin{equation}\label{eq:triplecost}
 C_{T_1T_2T_3}(x,y,z)=r_1(x)+r_2(y)+r_3(z)
       +\frac2{11}\{q(x+y)+q(y+z)\}+\frac15q(x+y+z),
\end{equation}
using $r_L^*$ for every L and the actual pressure residual
$r_{S,p}$ for each S. The four completed assertions are
\begin{center}\small
\begin{tabular}{@{}lll@{}}
\toprule
Labels & Domains by label & Lower bound\\
\midrule
$LLL$ & $L=[5/3,\infty)$ & $\lambda=9/100000$\\
$SLL_{\mathrm i}$ & $S_{\mathrm i},L,L$ & $\mu=293723/2400000000$\\
$LSL_{\mathrm i}$ & $L,S_{\mathrm i},L$ & $\mu=293723/2400000000$\\
$SLL_{\mathrm e}$ & $S_{\mathrm e},L,L$ & $\mu_e=28447/480000000$\\
\bottomrule
\end{tabular}
\end{center}
The short domains are those of \eqref{eq:enlargeddomains}. Full
reversal gives the corresponding $LLS$ assertions. Exchanging only
the two long arguments of $SLL$ is not a symmetry of
\eqref{eq:triplecost}; the two adjacent-sum slots would change.

Every residual and every remaining energy term in
\eqref{eq:triplecost} is nonnegative on its domain. This permits a
single variable's reserve to pay the whole target in a tail. For
the first three rows, write $\eta$ for their target. The long core is
\begin{equation}\label{eq:newlongcore}
 [17/10,\ \tL+\eta/\pe].
\end{equation}
Below its lower endpoint, $d\geq33/100$ and
$P_L(d)\geq P_L(33/100)>\max(\lambda,\mu)$ by exact rational
comparison. Above its upper endpoint, $\phi_L(x)\geq\eta$.
For an interior S, only the core
\begin{equation}\label{eq:newshortcore}
 [47/50,633/500]
\end{equation}
remains: for $s\leq47/50$, the positive polynomial in
\eqref{eq:shortreserve} at $d=3/25$ already exceeds $\mu$.
The upper endpoint is the enlarged gate, not an extrapolated
derivative range.

The inherited completed endpoint mixed cover uses the cores
\begin{equation}\label{eq:endpointcores}
 S:[103/100,317/250],\qquad
 L:[189/100,203/100+28447/81100].
\end{equation}
The long lower tail is paid because
$P_L(14/100)=78495/1200000000>\mu_e$.
For the short lower tail, $d\geq3/100$ and the last two terms of
\eqref{eq:shortreserve}, with $p=\pe$, already exceed $\mu_e$.
The long upper endpoint is precisely $\tL+\mu_e/\pe$.
All these cores and tails are closed; any overlap is harmless, and
there is no omitted boundary between them.

\section{Exact adaptive and atomic triple verification}

\subsection{An integer lower bound for whole boxes}

The triple covers reuse the completed grid-500 tables. Round each
finite rational core outwards to inclusive cell indices and take
their Cartesian product. For an adaptive box $[a_j,b_j]$, use the
full sum ranges in \eqref{eq:sumindices}, now for one, two, or
three coordinates. A range minimum supplies the lower $q$ value
over every required distance.

The interval and LL-replacement tests can be implemented with integer
fixed-point arithmetic of scale $S=2^{80}$. Set
$T_i=\floor{S Q_i}$, regarding the binary64 value $Q_i$ as its
exact rational value. For a single coordinate cell a safe residual
integer is
\begin{equation}\label{eq:fixedresidual}
 \max\left(0,\floor{T_i/6}
                    +\floor{S\phi(i/500)}\right),
\end{equation}
where $\phi$ is its actual short pressure residual or the
piecewise long pressure in \eqref{eq:worstlong}. Precompute this
residual bound for every coordinate cell. For a wider coordinate
interval, the implemented bound takes the minimum of those residual
bounds over all its cells; each individual cell bound is valid on
its entire closed cell. One could instead take a range minimum of
$T_i$ and use the leftmost pressure value, but that generally gives
a weaker lower bound and is not the recorded residual-table rule.
Flooring a negative pressure contribution means rounding toward
$-\infty$, not truncating toward zero. Each of the two lag-two
contributions and the lag-three contribution is also multiplied by
its exact coefficient and floored downwards. Accept only if the
sum of these integers is at least $\ceil{S\eta}$.

If two adjacent labels are L, one alternative bound replaces their
two residuals and connecting lag-two term by
$\floor{S K_L}$, retaining the other residual and unused energy
slots. When both adjacent pairs are LL, take the maximum of the
two alternatives, not their sum. The ordinary interval lower bound
is another alternative. Every accepted lower bound therefore uses
each of the six slots of \eqref{eq:triplecost} at most once.

Bisect the widest nonsingleton integer coordinate range as in
Section 4. Reversal can remove only boxes whose entire reflected
box is represented by the retained half: for $LLL$ and $LSL$ the
first and third label domains agree, whereas no internal reversal
reduction is made in $SLL$. A valid implementation retains
diagonal or overlapping ranges for further examination instead of
discarding them on the basis of a center ordering.

\subsection{A rigorous derivative from three exact point values}

The preceding tests leave some atomic cells
unresolved. The following exact quadratic test closes them without
requiring a finer scalar grid. Let $h=10^{-12}$. By
\eqref{eq:globalderivative}, $q'$ is Lipschitz with constant
$4\pi^2$. The centered difference is its average over
$[x-h,x+h]$, so
\begin{align}
 \left|q'(x)-\frac{q(x+h)-q(x-h)}{2h}\right|
 &\leq\frac{4\pi^2}{2h}\int_{-h}^h|v|\,dv
   =2\pi^2h.\label{eq:finiteerror}
\end{align}
Using the rational point evaluator for $q(x-h)$ and $q(x+h)$,
and $2\pi_{\hi}^2h$ as the error allowance, gives an exact
rational interval for $q'(x)$. Cancellation may widen this
interval, but cannot invalidate it. All required centers and
their two probes are above $.9$. The guarded below-one formula
in Section 3 is therefore available when a short midpoint is
below one; an $x\geq1$ routine is not silently applied outside
its domain.

\subsection{The full atomic Hessian and exact LDL test}

The long pressure kink is at $\tL=2.03$, exactly grid boundary
1015. Every closed atomic grid-500 cell lies on one side. Its
appropriate affine branch agrees with the true pressure at every
point of the closed cell, including a shared endpoint. The short
pressure is already affine. Thus on an atomic triple cell, with
coordinates $g=(x,y,z)$, the true cost has affine pressure and
Hessian
\begin{equation}\label{eq:triplehessian}
 \nabla^2 C(g)=\sum_{e=1}^6 c_e q''(a_e^Tg)a_ea_e^T,
\end{equation}
where the six pairs $(a_e,c_e)$ are
\begin{equation}\label{eq:incidences}
 \begin{gathered}
 ((1,0,0),1/6),\quad ((0,1,0),1/6),\quad((0,0,1),1/6),\\
 ((1,1,0),2/11),\quad((0,1,1),2/11),\quad((1,1,1),1/5).
 \end{gathered}
\end{equation}
For an atomic tuple $(i,j,k)$, the lag-two curvature ranges
include both indices $i+j,i+j+1$ and $j+k,j+k+1$. The
lag-three range includes $i+j+k,i+j+k+1,i+j+k+2$.
Let $\ell_e$ be the minimum of the corresponding complete
curvature range, interpreted as an exact rational lower bound.
Then
\begin{equation}\label{eq:tripleM}
 M=\sum_ec_e\ell_ea_ea_e^T\preceq\nabla^2C(g)
\end{equation}
throughout the atomic cell. Missing or nonfinite curvature disables
this test; it is never changed to zero to obtain positive pivots.

Let $g_0=((2i+1),(2j+1),(2k+1))/(2G)$ be the rational center.
Compute the LDL factorization $M=LDL^T$ with exact rational
arithmetic. If any pivot is nonpositive, this criterion gives no
acceptance. When every pivot is strictly positive, Taylor's
integral remainder on the segment inside the box gives
\begin{align}
 C(g)&\geq C(g_0)+\nabla C(g_0)^T(g-g_0)
                        +\tfrac12(g-g_0)^TM(g-g_0)\nonumber\\
 &\geq C(g_0)-\tfrac12\nabla C(g_0)^TM^{-1}\nabla C(g_0).
 \label{eq:atomictaylor}
\end{align}
The second line minimizes the lower quadratic over all of
$\R^3$. Neither a stationary point of $C$ nor the minimizing
point of that quadratic needs to be inside the cell.

Enclose $C(g_0)$ and its gradient using the exact point values
and \eqref{eq:finiteerror}. These evaluations use the actual
unclipped residual functions in \eqref{eq:triplecost}, not the
maximum with zero used in an independent interval lower bound.
For a gradient interval vector $G$, interval forward substitution
gives $Y=L^{-1}G$. If $Y_j=[l_j,u_j]$, put
\begin{equation}\label{eq:atomicU}
 U=\sum_{j=1}^3\frac{\max(|l_j|,|u_j|)^2}{D_j}.
\end{equation}
All $D_j$ are positive, so $U$ is an upper bound for the
quadratic debit despite dependencies between interval components.
The exact rational acceptance condition is
\begin{equation}\label{eq:atomicaccept}
 C(g_0)_{\lo}-U/2\geq\eta.
\end{equation}
No LL matching credit is added to this bound: the six original
slots have already been used to define $C$. Failed pivots,
inadequate lower bounds, domain-guard failures, or interrupted
point evaluation produce no acceptance.

The adaptive cover is complete only when both its pending stack
and its unresolved terminal list are empty. If a budget interrupt
occurs during atomic work, the current box is restored to pending
work. Node identities and exact covered atomic volumes are checked
in addition to, not instead of, the source's exhaustive traversal.

\section{Completed triple evidence and the exact assembly}

\subsection{What was actually executed}

The completed endpoint mixed computation needed no atomic Taylor
test. It visited 77 nodes with 38 splits and 39 accepted leaves:
37 interval leaves and 2 LL-replacement leaves. Ordinary and
optimized executions completed the full outward core in
\eqref{eq:endpointcores}.

The other three rows were completed with the atomic rule above,
both in ordinary Python and with \texttt{-OO}. The recorded
counts agree in both modes:
\begin{center}\small
\begin{tabular}{@{}lrrrrrr@{}}
\toprule
Case & Nodes & Splits & Interval & LL & Taylor & Reversal\\
\midrule
$LLL$ & 4117 & 2058 & 1059 & 9 & 868 & 123\\
$SLL_{\mathrm i}$ & 2553 & 1276 & 1106 & 2 & 169 & 0\\
$LSL_{\mathrm i}$ & 263 & 131 & 119 & 0 & 0 & 13\\
\bottomrule
\end{tabular}
\end{center}
Every row has one initial root and zero pending or unresolved
terminal boxes. The maximum depths are respectively 27, 27, and
22. The outward core atomic volumes are respectively
$80621568$, $45441792$, and $45441792$. They are volumes of
adaptive leaves, not counts of point samples. In the symmetric
rows, the proved and reversal-accounted volumes are respectively
$78952385+1669183$ for $LLL$ and
$42186303+3255489$ for $LSL_{\mathrm i}$; the $SLL_{\mathrm i}$
cover uses no reversal reduction.

These are completed executions of the distributed sources. All
1037 Taylor-accepted atomic cells were checked with rational matrix
lower bounds, LDL factors, center and gradient intervals, and
quadratic debits. The compact receipts retain exact canonical
digests and counts of these generated objects, not the detailed
matrices or point cache. The optional \texttt{--detailed} switch
saves those objects in a fresh execution. The ordinary and optimized
receipts agree after removal only of the optimization flag and
the three declared timing fields. Checking the Taylor cells alone
would still not prove an entire cover: the interval leaves, LL
leaves, legitimate symmetries, roots, splits and zero pending volume
remain indispensable.

\subsection{The 78-term allocation}

Partition the twelve gaps into the four consecutive triples
\begin{equation}\label{eq:partition}
 (g_1,g_2,g_3),\quad(g_4,g_5,g_6),\quad
 (g_7,g_8,g_9),\quad(g_{10},g_{11},g_{12}).
\end{equation}
Let $C_b^{\rm actual}$ denote \eqref{eq:triplecost} with the
actual pressure at every coordinate, rather than the dominated
long choice $r_L^*$. With $A_{\rm pat}$ from
\eqref{eq:patternAD}, direct allocation of terms gives
\begin{equation}\label{eq:assemblyidentity}
 F=A_{\rm pat}+\sum_{b=1}^4C_b^{\rm actual}+R,
 \qquad R\geq0.
\end{equation}
The allocation can be checked by the following inventory:
\begin{center}\small
\begin{tabular}{@{}lrrr@{}}
\toprule
Distance in gaps & Original slots & Four triples & Remainder\\
\midrule
$h=1$ & 12 & 12 & 0\\
$h=2$ & 11 & 8 & 3\\
$h=3$ & 10 & 4 & 6\\
$4\leq h\leq12$ & 45 & 0 & 45\\
\midrule
Total & 78 & 24 & 54\\
\bottomrule
\end{tabular}
\end{center}
Each used slot has its unchanged original coefficient, namely
$1/6$, $2/11$, or $1/5$. Each pressure occurs once, through its
baseline plus its residual. All 54 unused energy slots are
nonnegative. By \eqref{eq:worstlong}, each certified worst-pressure
triple is no greater than $C_b^{\rm actual}$.

If there is exactly one interior S, its baseline deficit is
\begin{equation}\label{eq:interiordeficit}
 \eps-A_{\rm pat}=\frac{1083}{8000000}
                         +\frac{97}{100}\pint
    =\frac{941723}{2400000000}=3\lambda+\mu.
\end{equation}
The S is the middle coordinate of its block at positions
$2,5,8,11$, using $LSL_{\mathrm i}$. At positions
$3,4,6,7,9,10$ it is an end coordinate, using $SLL_{\mathrm i}$
or its reversed $LLS_{\mathrm i}$ assertion. The other three
blocks are $LLL$. Thus both interior placements are proved, and
\begin{equation}\label{eq:interiorconclusion}
 F\geq\frac{31}{5000}
 \quad\hbox{on every retained interior-one-short region}.
\end{equation}
No additional positive uniform margin is claimed here.

For either endpoint S, the mixed bound is $\mu_e$. Its deficit
is
\begin{equation}\label{eq:endpointdeficit}
 \eps-A_{\rm pat}=\frac{143647}{480000000}
          =3\left(\frac1{12500}\right)+\mu_e.
\end{equation}
Since $\lambda-1/12500=1/100000$, the three all-long blocks
now improve this to
\begin{equation}\label{eq:endpointconclusion}
 F\geq0.00623
 \quad\hbox{on both retained endpoint-one-short regions}.
\end{equation}
If all twelve gaps are L, all four triples use $\lambda$ and
\begin{equation}\label{eq:alllongconclusion}
 F\geq\tL P+4\lambda
       =\frac{51397}{8000000}=0.006424625.
\end{equation}
There are thirteen labeled regions here: the all-long region and
twelve one-short regions, seven classes under full reversal. This
count is not a percentage of the full orthant that has been proved.

\section{Why the mixed credit was rebalanced, and what remains open}

The stronger standalone interior mixed credit
$\mu_{\rm old}=365723/2400000000$ is false. This can be
demonstrated without an optimizer or a cover: evaluate the exact
rational triple
\begin{equation}\label{eq:oldwitness}
 (x,y,z)=(1.05423,\ 2.009397,\ 2.026457).
\end{equation}
The point evaluator of Section 3 and the exact cost
\eqref{eq:triplecost} give
\begin{equation}\label{eq:oldwitnessvalue}
 C_{SLL_{\mathrm i}}(x,y,z)\in
 [0.000132480370716898210350,\ 0.000132480370716898210351].
\end{equation}
The upper endpoint is below $\mu_{\rm old}$ and above the new
target $\mu$. Both long coordinates are below $\tL$, so their
worst-pressure residuals equal the actual interior-pressure
residuals. The failure is therefore not an artifact of incompatible
pressure choices. It refutes that local credit only, not the full
objective \eqref{eq:F13}.

The successful correction has a precise budget interpretation.
Increasing each of the three $LLL$ credits in a one-short
decomposition from $.00008$ to $.00009$ lowers the required
mixed credit by $.00003$. That is exactly
$\mu_{\rm old}-\mu$. The smaller mixed credit is established
on its whole domain by the completed computations, not inferred
from the single point \eqref{eq:oldwitness}.

Combining the scalar gate with
\eqref{eq:interiorconclusion}--\eqref{eq:alllongconclusion}
now gives the precise remaining boundary:
\begin{equation}\label{eq:remaining}
 F(g)<31/5000\quad\Longrightarrow\quad
 \text{at least two coordinates have a retained short label}.
\end{equation}
Whether such a strict counterexample exists is unresolved here.
The multi-short regions have not been certified. No universal
thirteen-point floor, no extra global zero-count decimal, and no
worldwide record-priority assertion follows from this local
completion. In particular, the two seven-point inputs to the main
manuscript remain exactly \eqref{eq:rowC} and \eqref{eq:rowH}.

\section{Reproduction and the scope of the package}

A self-contained reproduction has two logically distinct checks.
The first evaluates the main manuscript's scalar count expressions
from the two exact row configurations. Such an evaluation is
conditional on the local rows; it is not a replay of them. The
second runs a complete local verifier with the actual profile,
pressure, floor, whole-cell tables, acceptance rules, and exhaustive
coverage. The distributed package identifies these roles separately
as \texttt{verify\_counts.py} and
\texttt{verify\_local\_rows.py}. Its README and manifest specify
the commands actually supplied, their optional supplementary
stages, and the execution status of each distributed source.

For the thirteen-point local results, the mathematical reproduction
sequence is: the rational scalar bootstrap and closed gate; the
short and long point checks together with the analytic derivative
proofs; generation or full revalidation of the six-mode scalar
tables; the LL square; the endpoint mixed cover; the three final
triple covers with their atomic rule; and the exact assembly.
The distributed \texttt{build\_supplement\_tables.py} constructs
only the scalar tables. The entry point
\texttt{verify\_supplement.py} separates the cases
\texttt{baselines}, \texttt{LL}, \texttt{LLL},
\texttt{SLL\_i}, \texttt{LSL\_i}, \texttt{SLL\_e}, and
\texttt{witness}; none of these is a full twelve-gap cover.
The scalar-gate command and the completed receipt filenames are
specified in the accompanying README. The program
\path{assemble_supplement.py} binds the gate and all four
source-generated triple receipts to the distributed sources and
scalar data, then checks the exact regional allocation. It is a
consistency-and-assembly checker, not an independent validator of
arbitrary supplied proof logs. The reconstruction commands establish
its local premises.
Every stage is described above. A finite table can be distributed as an acceleration, but a
verifier must know how to establish each of its whole-cell
enclosures. Simply matching a stored digest establishes only that
the same table has been read.

All acceptance guards must remain active under optimized Python;
they must use explicit checked conditions rather than removable
assertions. A bounded run may finish with an honest incomplete
status. A modified program needs its own source-bound
completed receipt before its output is described as a fresh replay.

This distinction is especially important here. The available
completed records establish the two adopted seven-point rows and
the specific thirteen-point local regions above, subject to the
explicit computational trust boundary. They do not establish a
complete thirteen-point target. The supplement supplies a linear,
independently readable proof of what those records mean, while
leaving that final unresolved question visible.

\begin{thebibliography}{9}
\bibitem{ainta}
Ainta, \emph{zeta-simple-zeros}, MIT-licensed source code,
commit \texttt{040c5e899e658aed7b56a2a87f501798fe10761d},
\url{https://github.com/ainta/zeta-simple-zeros/tree/040c5e899e658aed7b56a2a87f501798fe10761d}.
The original scalar-cell and convex-tangent engine and the cosine
seven-point verification are attributed to this source; adapted
portions retain its license.
\end{thebibliography}
\end{document}
